\documentclass[11pt]{article}

\usepackage{../ppackage}
\usepackage{bbm}
\usepackage{import}






\usepackage{tikz}
\usetikzlibrary{arrows.meta,patterns}

\usepackage{../../../Notes/tikzit}
\usepackage{../../../Notes/utf8math}

\input{../../../Notes/TIKZ/digraph.tikzstyles}



\usepackage{url}


\newcommand{\solution}{
\bigskip\noindent
	\textbf{Solution: \\}}
	


\DeclareMathOperator{\size}{size}
\DeclareMathOperator{\conv}{conv}
\newcommand{\SV}{\mathrm{SV}}
\newcommand{\bigO}{O}
\newcommand{\cut}{\mathrm{cut}}
\newcommand{\LLL}{\mathrm{LLL}}
\newcommand{\setR}{\mathbb{R}}
\newcommand{\setZ}{\mathbb{Z}}
\newcommand{\setQ}{\mathbb{Q}}
\newcommand{\setC}{\mathbb{C}}
\newcommand{\setN}{\mathbb{N}}
\newcommand{\wt}[1]{\widetilde{#1}}
\newcommand{\opt}{{\sc 0/1-opt}\xspace}
\newcommand{\aug}{{\sc 0/1-aug}\xspace}
\newcommand{\psep}{{\sc 0/1-psep}\xspace}
\newcommand{\sep}{{\sc 0/1-sep}\xspace}
\newcommand{\fopt}{{\sc 0/1-testopt\xspace} }

\newcommand{\hpp}{\mathrm{HPP}}
\newcommand{\nodes}{\mathcal{V}}
\newcommand{\vol}{\mathrm{vol}}
\newcommand{\diag}{\mathrm{diag}}
\newcommand{\arcs}{\mathcal{A}}
\newcommand{\edges}{\mathcal{E}}
\newcommand{\paths}{\mathscr{P}}
\newcommand{\cycles}{\mathcal{C}}




\newcommand{\K}{{\mathcal K}}
\newcommand{\A}{{A}}
\newcommand{\B}{{B}}
\newcommand{\T}{\mathscr{T}}
\newcommand{\eE}{\mathscr{E}}
\newcommand{\eS}{\mathscr{S}}
\newcommand{\eP}{\mathscr{P}}
\newcommand{\eM}{\mathscr{M}}



\newcommand{\transp}{^{\mathrm{T}}}

\newcommand{\smallmat}[1]{\left( \begin{smallmatrix} #1 \end{smallmatrix}\right)}

\newcommand{\mat}[1]{ \begin{pmatrix} #1 \end{pmatrix}}
\newcommand{\smat}[1]{ \big(\begin{smallmatrix} #1 \end{smallmatrix}\big)}

\newcommand{\pc}{\mathscr{P}}
\newcommand{\ob}{\mathscr{O}}
\newcommand{\odds}{\mathscr{W}}
\newcommand{\up}{\mathscr{U}}
\newcommand{\ef}{\mathscr{F}}
\newcommand{\eh}{\mathscr{H}}
\newcommand{\ev}{\mathscr{V}}
\newcommand{\ec}{\mathscr{C}}
\newcommand{\eu}{\mathscr{U}}

\newcommand{\lex}{\mathrm{lex}}

\renewcommand{\leq}{\leqslant}
\renewcommand{\geq}{\geqslant}









\newcommand{\linhull}{\mathrm{lin.hull}}
\newcommand{\affhull}{\mathrm{affine.hull}}
\newcommand{\charcone}{\mathrm{char.cone}}
\newcommand{\cone}{\mathrm{cone}}
\newcommand{\rank}{\mathrm{rank}}
\newcommand{\wb}[1]{\overline{#1}}



\usepackage{enumerate}

      
\institute{\'Ecole Polytechnique F\'ed\'erale de Lausanne}
\lecture{Discrete Optimization}
\faculty{Prof. Friedrich Eisenbrand}
\term{Spring 2026}
\publishdate{May 26, 2026}
\duedate{ }
\problemset{Problem Set~13}

\begin{document}
\makeheader

\begin{enumerate}[1)]


\item \label{arbitrage} Let \(C\) be a set of currencies and let \(r_{ij}\) denote the exchange
rate from currency \(i\) to currency \(j\). We say that there is an
\emph{arbitrage opportunity} if there exists a sequence 
\[
i_0,i_1,\ldots,i_k=i_0 ∈ C
\]
such that
\[
\prod_{\ell=0}^{k-1} r_{i_\ell i_{\ell+1}}>1.
\]

 Explain how Bellman--Ford can be used
 to detect such an opportunity.

\item \label{Floyd} (Floyd-Warshall algorithm) Let $G = (V,A)$ be a weighted directed graph with weights $c: A ⟶ ℝ$.
Define $D^{(k)}_{ij}$ to be the length of a shortest path from \(i\) to \(j\) whose \emph{internal}
vertices are only allowed to belong to $\{1,\dots,k\}$.
For \(k=0\), no internal vertices are allowed. Therefore
\[
D^{(0)}_{ij}
=
\begin{cases}
0, & i=j,\\
c_{ij}, & ij\in E,\\
+\infty, & \text{otherwise.}
\end{cases}
\]

\begin{enumerate}[i)] 
\item
 Suppose that $G$ does not contain a negative weight cycle. 
  Show the recurrence
  \[
    D^{(k)}_{ij}
    =
    \min\left\{
      D^{(k-1)}_{ij},
      D^{(k-1)}_{ik}+D^{(k-1)}_{kj}
    \right\}.
  \]
\item Prove that the matrix $D^{(n)}$ can be computed in time $O(n^3)$ steps.
\item  Show that $G$ contains a negative cycle if and only if there exists an $i ∈V$ such that $D^{(n)}_{ii} <0$. 
\end{enumerate}
\item \label{hamilton} (Hamilton path problem) The following problem is NP-complete. So far, no polynomial-time algorithm has been found to solve it:

  Given a directed graph $G = (V,A)$ and two vertices $s≠t ∈ V$, decide whether there exists a path from $s$ to $t$ of length $|V|-1$. This is a path that spans all the vertices of $G$.

  \begin{enumerate}[i)] 
  \item  Show that this problem can be solved in polynomial time, if the graph does not have any directed cycles.
  \item Conclude that the shortest path problem (under presence of negative  cycles) should be difficult. 
  \end{enumerate}

\item Consider the following problem. We are given $B ∈\mathbb{N}$, and a set of integer points $S= \{p ∈\setZ^n :
0 ≤p_i ≤B,\; ∀ i = 1,...,n\}$, whose points are all colored blue but one, which is red. We have an
oracle that, given vectors $l,r ∈\setR^n$, tells us whether the red point in $S$ is contained in the box
$S∩\{x ∈\setR^n : l_i ≤x_i ≤r_i,\;∀i = 1,...,n\}$ or not. Give an algorithm to find the red point using
$O(n \log(B))$ many oracle calls.


\end{enumerate}




  

\end{document}

%%% Local Variables:
%%% mode: latex
%%% TeX-master: t
%%% End:
